\section{Pseudometric spaces}

\import{surreal/real-field.tex}

Pseudometrics take values in the canonical surreal real field.  Distinct
points may have distance zero; metrics will add precisely the missing
separation law.

\begin{struct}\label{pseudometric_space}
    A pseudometric space $M$ is a onesorted structure equipped with
    \begin{enumerate}
        \item\label{pseudometric_distance_graph}
            a distance map
            $\metricDistance[M]:
                \carrier[M]\times\carrier[M]\to\reals$;
    \end{enumerate}
    such that
    \begin{enumerate}
        \item\label{pseudometric_distance_nonnegative}
            for all $x,y\in\carrier[M]$ we have
                $\surrealLe{\surrealZero}{\metricDistance[M](x,y)}$.
        \item\label{pseudometric_distance_self}
            for every $x\in\carrier[M]$ we have
                $\metricDistance[M](x,x)=\surrealZero$.
        \item\label{pseudometric_distance_symmetric}
            for all $x,y\in\carrier[M]$ we have
                $\metricDistance[M](x,y)=\metricDistance[M](y,x)$.
        \item\label{pseudometric_triangle}
            for all $x,y,z\in\carrier[M]$ we have
                $\surrealLe{\metricDistance[M](x,z)}
                    {\surrealAdd{\metricDistance[M](x,y)}
                        {\metricDistance[M](y,z)}}$.
    \end{enumerate}
\end{struct}

\begin{proposition}\label{pseudometric_distance_type}
    Let $M$ be a pseudometric space.
    Let $x,y\in\carrier[M]$.
    Then $\metricDistance[M](x,y)\in\reals$.
\end{proposition}
\begin{proof}
    $(x,y)\in\carrier[M]\times\carrier[M]$
        by \obv{times_tuple_intro}.
    Follows by \obv{pseudometric_distance_graph,funs_type_apply}.
\end{proof}
